\section{A uniform positive ionization gap}\label{a:sec:gap}

The energy bound of Section~\ref{a:sec:neutral} makes a distant insertion
trial effective in every singly ionized atom.  Averaging the insertion
center recovers its unit unscreened charge; a complementary localized
state controls the attraction lost when a hole is made for the orbital.

\begin{proposition}[Positive ionization energy]\label{a:prop:tail-gap}
There are universal constants $c>0$ and $Z_2<\infty$ such that
\begin{equation}\label{a:eq:tail-gap}
 E_{Z-1}-E_Z\ge c\qquad(Z\ge Z_2).
\end{equation}
\end{proposition}

\begin{proof}
For $Z\ge\max(Z_0,2)$, take a normalized ground state $\Phi$ of
$H_{Z-1,Z}$.  It exists by Lemma~\ref{lem:atomic-binding}, and its offset
from $E$ is bounded by Proposition~\ref{a:prop:tail-offset}.  In this proof
$S\ge1$ will be a large universal constant, fixed at the end, and
$b=S^{3/5}$.  Angle brackets denote normalized Lebesgue averaging over
the centers $\mathcal A_S=\{y:S\le|y|\le2S\}$.

Use the real normalized radial packet $g_b$ from
Section~\ref{sec:localization}.  Choose smooth real functions
$c_y,o_y$ with $c_y^2+o_y^2=1$, $c_y=0$ on $B(y,2b)$, $c_y=1$
outside $B(y,3b)$, and
\[
 |\nabla c_y|^2+|\nabla o_y|^2
 \le Cb^{-2}\1_{B(y,3b)}.
\]
Set
\[
 F_y=\prod_{i=1}^{Z-1}c_y(x_i),\quad
 n_y^{\rm near}=\#\{i:|x_i-y|<3b\},\quad
 N_y=\E_\Phi n_y^{\rm near},\quad
 A_y=\E_\Phi F_y^2.
\]
The unnormalized core $F_y\Phi$ has no particles in $B(y,2b)$.
Wedge this core with the normalized one-spin orbital $g_b(\cdot-y)$;
the resulting unnormalized $Z$-electron vector has squared norm $A_y$.
Spatial separation makes the
terms assigning different particle labels to the two groups orthogonal.
It also makes its energy equal to the core energy, plus the orbital
energy times $A_y$, plus the direct cross repulsion.  In particular,
with the raw-configuration field
\[
 H_y^{\rm add}
 =\int g_b(x-y)^2
       \left(V(x)-\sum_{i=1}^{Z-1}K(x-x_i)\right)\,\dd x,
\]
the multiplier identity \eqref{eq:sector-multiplier} and the variational
principle give
\begin{align}
 (E_Z-E_{Z-1})A_y
 &\le \frac12\E_\Phi|\nabla F_y|^2
       +\frac{\|\nabla g\|_2^2}{2b^2}A_y
       -\E_\Phi[F_y^2H_y^{\rm add}]\notag\\
 &\le Cb^{-2}(1+N_y)-\E_\Phi[F_y^2H_y^{\rm add}].
 \label{a:eq:tail-insertion}
\end{align}
This inequality also holds if $A_y=0$, since its derivation uses the
unnormalized trial and its zero norm.

For every raw configuration, averaging the full insertion field gives
\begin{equation}\label{a:eq:tail-unit-charge}
 \langle H_y^{\rm add}\rangle\ge\frac1{2S}.
\end{equation}
To verify this without a symmetry assumption on $\Phi$, fix the center
radius $r\in[S,2S]$.  The nuclear packet average is $Z/r$ since $b<r$.
For any fixed electron position $z$, the spherical average in $y$ of its
packet potential is
\[
 \int g_b(v)^2\frac1{\max(r,|z-v|)}\,\dd v\le\frac1r.
\]
This follows by first averaging $K(y+v-z)$ over $|y|=r$ and then
integrating $v$.  There are $Z-1$ electrons, so the averaged field is at
least $1/r\ge1/(2S)$.  Radial averaging proves
\eqref{a:eq:tail-unit-charge}.

The mean number of particles near a proposed center is small.  For
large $S$, $3b<S/2$, and a particle can lie in $B(y,3b)$ for some
$y\in\mathcal A_S$ only if $S/2\le|x_i|\le3S$.  The fraction of such
centers for any particle is at most $Cb^3/S^3$.  Fubini and
\eqref{a:eq:screen-annular}, with the bounded offset of $\Phi$, yield
\begin{equation}\label{a:eq:tail-insertion-overlap}
 \langle N_y\rangle
 \le C(b/S)^3\E_\Phi\#\{i:S/2\le|x_i|\le3S\}
 \le C(b/S)^3(1+\sqrt S)
 \le CS^{-7/10}.
\end{equation}
Writing $p_y=1-A_y=\E_\Phi(1-F_y^2)$, the product inequality
$1-\prod_i c_y(x_i)^2\le\sum_i(1-c_y(x_i)^2)$ gives $p_y\le N_y$.

It remains to bound the potentially positive field lost on the
discarded configurations.  The symmetric multiplier
$G_y=\sqrt{1-F_y^2}$ is Lipschitz, even where it vanishes, and
\begin{equation}\label{a:eq:tail-complement-gradient}
 |\nabla G_y|^2\le Cb^{-2}n_y^{\rm near}.
\end{equation}
For a direct justification, take all products of $c_y$ and $o_y$ in
the particle variables except the all-core product.  Their Euclidean
norm is $G_y$.  The squared gradient of a norm is bounded by the sum
of the squared gradients of its components; the full product partition
has squared gradient sum
$\sum_i(|\nabla c_y|^2+|\nabla o_y|^2)(x_i)$.
This proves \eqref{a:eq:tail-complement-gradient} and the stated
Lipschitz property.

If $p_y>0$, put $\Phi_y^{\rm loss}=G_y\Phi/\sqrt{p_y}$.  This is a
normalized antisymmetric form-domain vector.  By
\eqref{eq:sector-multiplier}, \eqref{a:eq:tail-offset}, and
\eqref{a:eq:tail-complement-gradient}, its offset satisfies
\begin{equation}\label{a:eq:tail-complement-offset}
 \delta_y\le C+Cb^{-2}N_y/p_y.
\end{equation}
Choose $S$ large enough also that $b<A S/L$.  At each center $y$ in
the annulus, the nucleus and every electron included in $J_y$ lie
outside the radial packet.  Their packet averages equal their point
potentials at $y$.  All omitted electron terms enter with a negative
sign, whence $H_y^{\rm add}\le J_y$ pointwise in the configuration.
Using \eqref{a:eq:screen-local} with no initial cut, now in the state
$\Phi_y^{\rm loss}$, gives
\begin{align}
 \E_\Phi[(1-F_y^2)H_y^{\rm add}]
 &=p_y\E_{\Phi_y^{\rm loss}}H_y^{\rm add}\notag\\
 &\le Cp_y\big(S^{-1}+\sqrt{\delta_y}\,S^{-1/2}\big)\notag\\
 &\le C\left[
 p_y(S^{-1}+S^{-1/2})
 +b^{-1}S^{-1/2}\sqrt{p_yN_y}\right].
 \label{a:eq:tail-loss-bound}
\end{align}
Here we may assume $S/L\ge1$; $L$ is fixed.  If $p_y=0$, the left
side is zero and the same bound holds.  All expectations in this
calculation are finite for a fixed system by the form bounds (the
packet potential is bounded as well).

Since $p_y\le N_y$, \eqref{a:eq:tail-insertion-overlap} implies
\[
 \left\langle\E_\Phi[(1-F_y^2)H_y^{\rm add}]\right\rangle
 \le C\left(S^{-17/10}+S^{-6/5}+S^{-9/5}\right)
 =o(S^{-1}).
\]
Also
\[
 Cb^{-2}(1+\langle N_y\rangle)
 \le C(S^{-6/5}+S^{-19/10})=o(S^{-1}).
\]
These estimates explain the width choice: any $b=S^a$ with
$1/2<a<2/3$ would make the kinetic term and the largest averaged
field loss smaller than $S^{-1}$.
All these limits are uniform in $Z$ and in the ground state $\Phi$.
Average \eqref{a:eq:tail-insertion} and write
$F_y^2H_y^{\rm add}=H_y^{\rm add}-(1-F_y^2)H_y^{\rm add}$.
Equations~\eqref{a:eq:tail-unit-charge} and \eqref{a:eq:tail-loss-bound}
then show, for one sufficiently large universal choice of $S$,
\[
 (E_Z-E_{Z-1})\langle A_y\rangle\le-\frac1{4S}.
\]
In particular $\langle A_y\rangle>0$.  As it is at most one, dividing
gives $E_{Z-1}-E_Z\ge1/(4S)$.  This proves
\eqref{a:eq:tail-gap}.
\end{proof}

\begin{proof}[Proof of Corollary~\ref{cor:first-ionization}]
In the notation $E_Z(n)=E(n,Z)$ of the corollary,
Proposition~\ref{a:prop:tail-offset} gives universal $C_0,Z_0<\infty$
such that, for $Z\ge Z_0$,
\[
 I_1(Z)=E_Z(Z-1)-E_Z(Z)
 \le E_Z(Z-1)-\inf_{n\ge0}E_Z(n)\le C_0.
\]
The first inequality uses $\inf_{n\ge0}E_Z(n)\le E_Z(Z)$.
Proposition~\ref{a:prop:tail-gap} gives universal $c_0>0$ and
$Z_2<\infty$ such that $I_1(Z)\ge c_0$ for $Z\ge Z_2$.

Choose an integer $Z_*\ge\max\{2,Z_0,Z_2\}$. For each integer
$1\le Z<Z_*$, the finite-sector Coulomb forms are bounded below and admit
finite-energy trial states, so $E_Z(Z-1)$ and $E_Z(Z)$ are finite; the
vacuum energy is $E_Z(0)=0$. Lemma~\ref{lem:atomic-binding}, applied with
$n=Z$, gives $I_1(Z)>0$. Thus all the finitely many remaining gaps lie
in $(0,\infty)$. Set
\[
 c=\min\{c_0,I_1(1),\ldots,I_1(Z_*-1)\},
 \qquad
 C=\max\{C_0,I_1(1),\ldots,I_1(Z_*-1)\}.
\]
These constants satisfy $0<c\le C<\infty$. They are universal because
$Z_*$ is universal and the exceptional gaps belong to finitely many atoms
in the fixed normalization. The large-$Z$ bounds and these choices prove
the assertion for every integer $Z\ge1$.
\end{proof}
